%% ==========================================================================
%% DRAFT -- August 11, 2026. Author byline: Daniel Kirtchakov (05oz,
%% halfounce.io; no institutional affiliation; ORCID 0009-0009-5213-4098).
%%
%% Result referee-CONFIRMED 2026-08-11 (final adversarial pass, fresh code
%% throughout; the referee's verdict record and code are retained by the
%% author and are not deposited -- the deposit carries the review log and
%% the separate k(6,3) verdict record). Novelty swept clean
%% 2026-08-05 and 2026-08-11 (erdosproblems.com/112 fetched live: OPEN,
%% 0 comments, 0 claimed proofs, no values recorded; arXiv: nothing on
%% r(I_m,L_n) beyond IRW; IRW citing literature: one off-target citation).
%% Sweep re-run on the drafting date (2026-08-11) before this file was
%% written: unchanged.
%%
%% Framing constraints honored throughout (PROTOCOL section 11):
%%   * The claim is stated at exactly the referee-permitted strength:
%%     k(3,4) = r(I_3,L_4) = 21, determined by an explicit 20-vertex witness
%%     plus a 346-cube LRAT-certified exhaustion at N=21 with independently
%%     audited completeness. Previous published bounds: 9 <= k(3,4) <= 25
%%     (lower: monotonicity from Bermond's k(3,3)=9; upper: IRW 2021).
%%   * Secondary results at their permitted strength only: the N=22 and
%%     N=23 certified exhaustions are REDUNDANT confirmations; and
%%     29 <= k(6,3) <= 33 (lower bound ours, upper bound IRW's).
%%     k(6,3) remains OPEN and no more is claimed.
%%   * Computation files live in the repository, not the paper (one
%%     Certification paragraph, no file inventories, no chronology in the
%%     body). The note opens with the mathematical question and closes with
%%     the questions it opens.
%% ==========================================================================
\documentclass[11pt]{amsart}

\usepackage[margin=1in]{geometry}
\usepackage{amsmath,amssymb}
\usepackage{booktabs}
\usepackage{array}
\usepackage[colorlinks=true,allcolors=blue]{hyperref}

\emergencystretch=3em
\tolerance=1200
\hbadness=4000

\newcommand{\TT}[1]{\mathrm{TT}_{#1}}
\newcommand{\I}[1]{I_{#1}}
\newcommand{\Np}{N^{+}}
\newcommand{\Nm}{N^{-}}
\newcommand{\NI}{I}

\theoremstyle{plain}
\newtheorem{theorem}{Theorem}[section]
\newtheorem{lemma}[theorem]{Lemma}
\newtheorem{proposition}[theorem]{Proposition}
\newtheorem{corollary}[theorem]{Corollary}
\theoremstyle{definition}
\newtheorem{definition}[theorem]{Definition}
\theoremstyle{remark}
\newtheorem{remark}[theorem]{Remark}
\newtheorem{question}[theorem]{Question}

\begin{document}

\title[The oriented Ramsey number $k(3,4)=21$]{An Erd\H{o}s--Rado
oriented Ramsey number determined: $k(3,4)=r(I_3,L_4)=21$, by explicit
witness and certified exhaustion}

\author{Daniel Kirtchakov}
\address{Independent researcher (\texttt{05oz}); no institutional
affiliation. ORCID:
\href{https://orcid.org/0009-0009-5213-4098}{0009-0009-5213-4098}.}
\email{daniel@halfounce.io}
\urladdr{https://halfounce.io}

\thanks{\emph{Computation and authorship.} All encodings, decompositions,
searches, verifiers, and audits in this work were produced by
\textbf{Claude} (Anthropic), directed by the author, on a single
Apple~M4 laptop. The independent proof checker imports only the Python
standard library and shares no code with the solving pipeline; the
adversarial referee of \S\ref{ssec:referee} was a separate agent instance
given no access to the pipeline's code and wrote its own code throughout,
including its own LRAT checker. This is a factual methods statement, and
it is part of the point of the note: the artifacts are designed so that
the provenance of the \emph{search} is irrelevant to the validity of the
\emph{result}. External tools used by the pipeline: CaDiCaL~3.0.1
(\texttt{--lrat}) and Python~3.}

\thanks{\emph{Prior-art record.} The primary source \cite{IRW21} was read
in full on August~5, 2026 (arXiv v3; a copy is archived in the author's
working tree, not redistributed here, and section references follow its
numbering). The Erd\H{o}s
problems entry \cite{Blo} was fetched live
on August~5 and again on August~11, 2026 (including once more immediately
before this draft was written): status \textsc{open}, zero comments, zero
claimed proofs, no exact value of any $k(n,m)$ recorded. An arXiv sweep
of the same dates found no paper on $r(I_m,L_n)$ other than \cite{IRW21},
whose single recorded citation is off-target. The line has been dormant
since 2020/2021. \emph{Draft of August~11, 2026.}}

\subjclass[2020]{Primary 05C55; Secondary 05C20, 05D10, 68V15}
\keywords{Oriented graph, Ramsey number, transitive tournament,
independent set, Erd\H{o}s problem, SAT, cube-and-conquer, LRAT,
certified computation}

\date{August 11, 2026}

\begin{abstract}
Let $k(n,m)$ be the least $N$ such that every directed graph on $N$
vertices contains an independent set of size $n$ or a transitive
tournament on $m$ vertices. Determining $k(n,m)$ is a problem of
Erd\H{o}s and Rado (1967), \#112 in the Erd\H{o}s problems collection;
it coincides with the oriented-graph Ramsey number $r(I_n,L_m)$ studied
by Bermond and by Ihringer, Rajendraprasad, and Weinert (IRW). Beyond the classical tournament column $k(2,m)$, the known
exact values were $k(3,3)=9$ (Bermond 1974) and $k(4,3)=15$, $k(5,3)=23$
(IRW 2021); for $k(3,4)$, the case IRW single out as the next feasible
one, the published bounds were $9\le k(3,4)\le 25$. We determine
$k(3,4)=r(I_3,L_4)=21$.

The lower bound is an explicit oriented graph on $20$ vertices, printed
in full in \S\ref{sec:witness} and checkable in milliseconds. The upper
bound is an exhaustion of the $21$-vertex case: a neighbourhood
decomposition splits it into $346$ propositional instances; CaDiCaL
refuted every instance with an LRAT certificate, and every certificate
was replayed by an independently written standard-library checker. Completeness of the decomposition was
audited at content level: complete block inventories (three independent
enumerations), a base encoding with no symmetry breaking whose models
are exactly the $\{I_3,\TT{4}\}$-free oriented graphs (brute-forced at
small orders), and a one-to-one clause-multiset match of the $346$
instance files against the decomposition re-derived from the
definitions. With logically redundant confirming exhaustions at
$N=22,23,24$, the corpus is $445$ certificates and just over $10^9$
checked proof steps; certificates regenerate bit-for-bit from the
recorded solver commands ($11$ of $11$ attempts). Every component was then confirmed by an
adversarial referee writing fresh code throughout. From the same
campaign: $29\le k(6,3)\le 33$, the lower bound new, by an explicit
Cayley witness on $28$ vertices; the upper bound is IRW's. We close with
the questions the computation opens: the structure and possible
uniqueness of the $20$-vertex extremal graph, the $k(6,3)$ gap, the
feasibility of $k(4,4)$ (the new value gives $k(4,4)\le 50$), and what
a human-readable proof would require.
\end{abstract}

\maketitle

\section{Introduction}\label{sec:intro}

\subsection{The question}\label{ssec:question}
An \emph{oriented graph} is a loopless directed graph with at most one
arc between any two vertices; a \emph{tournament} is an oriented graph
in which every pair is adjacent; a tournament is \emph{transitive} when
its arc relation is a linear order. Write $\I{n}$ for the independent set
on $n$ vertices and $L_m$ (equivalently $\TT{m}$) for the transitive
tournament on $m$ vertices. Following \cite{IRW21}, $r(\I{n},L_m)$ is
the least $N$ such that every oriented graph on $N$ vertices contains an
independent set of size $n$ or a transitive tournament on $m$ vertices
(for tournaments, subdigraph and induced-subdigraph containment
coincide).

Erd\H{o}s and Rado \cite{ErRa67} asked for the same quantity over
directed graphs: let $k(n,m)$ be minimal such that any directed graph on
$k(n,m)$ vertices contains an independent set of size $n$ or a
transitive tournament of size $m$. The two functions are equal
(Remark~\ref{rem:bridge}), and determining $k(n,m)$ is \#112 in the
Erd\H{o}s problems collection \cite{Blo}, recorded there as open.
That entry records the following. Erd\H{o}s and Rado
proved $k(n,m)\ll_m n^{m-1}$, with the explicit bound
$k(n,m)\le\bigl(2^{m-1}(n-1)^m+n-2\bigr)/(2n-3)$; Larson and Mitchell
\cite{LaMi97} sharpened the dependence on $m$, in particular to
$k(n,3)\le n^2$ --- a bound also stated as
\cite[Lemma~2.4]{IRW21}, where it is attributed to \cite{LaMi97}; and
Z.~Hunter observed
$R(n,m)\le k(n,m)\le R(n,m,m)$, hence $k(n,m)\le 3^{n+2m}$.

The exact values previously known form two families. The tournament
column is classical: $k(2,3)=4$, $k(2,4)=8$, $k(2,5)=14$, $k(2,6)=28$
(see \cite{IRW21} and the references there). In the third-column
direction, Bermond \cite{Ber74} proved $k(3,3)=9$; the extremal graph on
$8$ vertices is unique up to isomorphism by
\cite[Lemma~3.1]{IRW21}, which also gives its description as the
circulant on $\mathbb{Z}_8$ with connection set $\{1,6\}$. Ihringer,
Rajendraprasad, and Weinert
\cite{IRW21} proved $k(4,3)=15$ and $k(5,3)=23$, gave the general bound
$k(m,3)\le m^2-m+3$, and --- combining their upper bound with Kim's
lower bound for $r(\I{m},K_3)$ --- determined the order of growth
$k(m,3)=\Theta(m^2/\log m)$. For $k(3,4)$ --- the case \cite{IRW21}
single out in closing as the next one within reach --- the record stood
at
\[
9 \;\le\; k(3,4) \;\le\; 25,
\]
the lower bound by monotonicity from $k(3,3)=9$, the upper bound from
the recursion of \cite[Lemma~2.3]{IRW21},
$r(\I3,L_4)\le 2\,r(\I3,L_3)+r(\I2,L_4)-1=25$; the same value $25$ is
what \cite[Proposition~6.1]{IRW21}, the small-parameter formula the
authors describe as the state of the art for small $m$ and $n$, returns
at $(m,n)=(3,4)$.

\subsection{Result}\label{ssec:result}

\begin{theorem}\label{thm:main}
$k(3,4)=r(I_3,L_4)=21$. That is: every directed graph --- equivalently,
every oriented graph --- on $21$ vertices contains an independent set of
size $3$ or a transitive tournament on $4$ vertices, and there is an
explicit oriented graph on $20$ vertices, exhibited in
\S\ref{sec:witness}, containing neither.
\end{theorem}

This determines the first previously unknown value in the $k(3,m)$
column of Erd\H{o}s Problem \#112. The proof is machine-checkable end to
end. The lower bound is the explicit witness of \S\ref{sec:witness}. The
upper bound is an exhaustive case decomposition of the $21$-vertex
problem into $346$ propositional instances (\S\ref{sec:upper}), each
refuted by CaDiCaL~3.0.1 \cite{CaDiCaL} with
an LRAT certificate \cite{LRAT}, in the spirit of
cube-and-conquer \cite{HKWB11} and of the certified exhaustions of
\cite{HKM16}; the completeness of the decomposition --- that the $346$
cases cover everything, with no symmetry breaking hidden in the
encoding --- was itself audited at content level
(\S\ref{sec:verif}). Every certificate was verified by an independently
written checker, and twelve of them by a second, unrelated one; the
whole was confirmed by a final
adversarial referee (\S\ref{ssec:referee}).

Two secondary results from the same campaign are recorded at exact
strength. The certified exhaustions at $N=22,23$ (and a structured
instance at $N=24$) are
logically redundant given Theorem~\ref{thm:main} but stand as
independent confirmations (\S\ref{ssec:redundant}). And:

\begin{proposition}\label{prop:k63}
$29\le k(6,3)\le 33$. The lower bound is new, by the explicit
$28$-vertex witness $\mathrm{Cay}(\mathbb{Z}_{28},\{3,8,10,12,17\})$;
the upper bound is $m^2-m+3$ at $m=6$, due to \cite{IRW21}.
\end{proposition}

\noindent $k(6,3)$ remains open, and nothing stronger about it is
claimed here.

\begin{remark}[Directed versus oriented graphs]\label{rem:bridge}
Erd\H{o}s Problem \#112 is phrased for directed graphs, which may carry
$2$-cycles; the computation here is over oriented graphs. The two Ramsey
functions coincide. Every oriented graph is a directed graph, so
$k(n,m)\ge r(\I{n},L_m)$. Conversely, a directed graph avoiding both
patterns yields an oriented one on the same vertex set: delete one arc
of each $2$-cycle. Deletion preserves adjacency and non-adjacency of
every pair --- hence preserves every independent set --- and only
removes transitive subtournaments; so $k(n,m)\le r(\I{n},L_m)$. The
extremal orders are equal, and the value $21$ answers the problem as
posed.
\end{remark}

\subsection{What is not claimed}\label{ssec:notclaimed}
We do not claim to resolve Erd\H{o}s Problem \#112. The problem asks for
$k(n,m)$ as a function of two parameters and is not settled by any
finite computation; a single new entry in its table leaves it open. We
do not claim a human-readable proof of the upper bound: what a
skeptic must believe and can replay is stated in \S\ref{sec:tcb}; what
a human proof would require, in
\S\ref{sec:questions}. We do not claim uniqueness of the $20$-vertex
extremal graph (Question~\ref{q:unique}). The asymptotics of $k(n,m)$
are untouched: the best general results remain those of
\cite{ErRa67,LaMi97,IRW21}. We do not claim
$k(6,3)$: Proposition~\ref{prop:k63} narrows its range and no more. The
prior bounds improved on here are credited at their sources: the
interval $[9,25]$ is \cite{Ber74} plus \cite{IRW21}.

\section{The lower bound: an explicit graph on $20$
vertices}\label{sec:witness}

Table~\ref{tab:witness} defines an oriented graph $W$ on vertex set
$\{0,\dots,19\}$ by its out-neighbourhoods; $W$ has $126$ arcs.

\begin{table}[ht]
\begin{center}
\small
\begin{tabular}{r@{\ $\to$\ }l@{\qquad}r@{\ $\to$\ }l}
\toprule
$0$ & $1,3,4,7,9,18$        & $10$ & $2,3,5,6,17,19$ \\
$1$ & $2,4,6,7,10,14,16$    & $11$ & $2,4,6,8,9,10$ \\
$2$ & $0,8,12,16,17,18,19$  & $12$ & $9,10,11,15,17,19$ \\
$3$ & $1,6,11,12,13,16$     & $13$ & $0,1,9,10,11,14,17$ \\
$4$ & $3,5,10,12,13,14$     & $14$ & $0,3,6,16,18,19$ \\
$5$ & $2,6,11,12,13,15,18$  & $15$ & $0,1,2,8,9,14$ \\
$6$ & $4,8,12,13,17,19$     & $16$ & $0,4,7,13,15,17$ \\
$7$ & $2,3,13,14,15,18$     & $17$ & $1,3,4,5,8,11$ \\
$8$ & $3,5,7,9,10,12,18$    & $18$ & $1,9,13,15,16,19$ \\
$9$ & $2,5,7,14,16,19$      & $19$ & $0,4,5,7,8,11,15$ \\
\bottomrule
\end{tabular}
\end{center}
\caption{The $20$-vertex witness $W$: out-neighbourhoods.}
\label{tab:witness}
\end{table}

\begin{theorem}\label{thm:witness}
$W$ contains no independent set of size $3$ and no transitive tournament
on $4$ vertices. Hence $k(3,4)\ge 21$.
\end{theorem}

The check is a finite verification over the $\binom{20}{3}=1140$ triples
and $\binom{20}{4}=4845$ quadruples, and runs in milliseconds; it was
performed by independently written verifiers using three distinct
transitivity criteria (among them the score-sequence criterion:
transitive if and only if the scores are $\{0,1,2,3\}$), and again from
scratch by the adversarial referee of \S\ref{ssec:referee}, who
additionally pushed a randomly relabelled copy of $W$ through the full
canonicalization of \S\ref{ssec:canon} at every one of its $20$
vertices.

Three facts about $W$, directly checkable from Table~\ref{tab:witness},
are worth recording for \S\ref{sec:questions}. Every out- and in-degree
is $6$ or $7$, and the vertex types
$\bigl(d^+(v),\,d^-(v),\,|\NI(v)|\bigr)$ --- out-degree, in-degree,
number of non-neighbours --- take exactly four values: $(6,6,7)$ at
eleven vertices, and $(6,7,6)$, $(7,6,6)$, $(7,7,5)$ at three vertices
each. In particular $W$ is not vertex-transitive. Finally, at each
vertex of type $(6,6,7)$ the non-neighbourhood is a $7$-vertex
$\TT{4}$-free tournament, hence (Proposition~\ref{prop:blocks}) a copy
of the Paley tournament $QR_7$.

\section{The upper bound: exhaustion at $N=21$}\label{sec:upper}

Call an oriented graph \emph{free} if it contains neither $\I3$ nor
$\TT{4}$. The upper bound is the assertion that no free graph on $21$
vertices exists. The exhaustion has three mathematical ingredients ---
a neighbourhood lemma, two vanishing endpoints, and complete
inventories of the neighbourhood blocks --- and one computational
engine: $346$ certified refutations.

\subsection{The case decomposition}\label{ssec:decomp}

For a vertex $w$ of an oriented graph $D$, write $\Np(w)$, $\Nm(w)$,
$\NI(w)$ for the sets of out-neighbours, in-neighbours, and
non-neighbours of $w$.

\begin{lemma}\label{lem:nbhd}
Let $D$ be free and $w$ a vertex of $D$. Then
(i) $\Np(w)$ and $\Nm(w)$ each induce an oriented graph with no $\I3$
and no $\TT{3}$;
(ii) $\NI(w)$ induces a tournament with no $\TT{4}$.
\end{lemma}

\begin{proof}
(i) An independent triple inside $\Np(w)$ is one in $D$. A transitive
triangle inside $\Np(w)$ extends by $w$, a common in-neighbour of its
three vertices, to a $\TT{4}$ in $D$. The argument for $\Nm(w)$ is the
same with $w$ as common out-neighbour (sink side).
(ii) Two non-adjacent vertices of $\NI(w)$ together with $w$ form an
$\I3$, so $\NI(w)$ induces a tournament; a $\TT{4}$ inside it is one in
$D$.
\end{proof}

\begin{lemma}\label{lem:endpoints}
There is no $\{\I3,\TT{3}\}$-free oriented graph on $9$ vertices, and
no $\TT{4}$-free tournament on $8$ vertices.
\end{lemma}

Lemma~\ref{lem:endpoints} is equivalent to $k(3,3)=9$ \cite{Ber74} and
to the classical $k(2,4)=8$, but the exhaustion does not rest on the
literature for it: both statements were established internally four
ways --- by direct enumeration of all labelled graphs of the two
families (the labelled counts of $\{\I3,\TT{3}\}$-free graphs on
$1,\dots,9$ vertices are $1$, $3$, $20$, $224$, $2554$, $18370$,
$30960$, $5040$, $0$; of $\TT{4}$-free tournaments on $1,\dots,8$
vertices, $1$, $2$, $8$, $40$, $184$, $240$, $240$, $0$), re-derived
twice more by independent enumerations during the audits, and by
LRAT-certified refutations of the corresponding instances, replayed by
both checkers.

By Lemma~\ref{lem:nbhd} and Lemma~\ref{lem:endpoints}, in a free graph
every vertex $w$ has $d^+(w)\le 8$, $d^-(w)\le 8$, $|\NI(w)|\le 7$;
counting $N=1+d^++d^-+|\NI|\le 24$ caps free graphs at $24$ vertices,
hence $k(3,4)\le 25$ by pure arithmetic. This is the specialization to
$(3,4)$ of the argument of \cite[Lemma~2.3]{IRW21}, and it is exactly
where the previous record stood. The exhaustion works at $N=21$: there
$(p,q,s):=(d^+(w),d^-(w),|\NI(w)|)$ satisfies $p+q+s=20$ with
$p,q\le 8$, $s\le 7$, which also forces $p,q\ge 5$ and $s\ge 4$.

\begin{proposition}[Block inventories]\label{prop:blocks}
Up to isomorphism, the $\{\I3,\TT{3}\}$-free oriented graphs on
$5,6,7,8$ vertices number $25$, $31$, $7$, $1$ respectively --- the
unique one on $8$ vertices being Bermond's circulant, in agreement with
\cite[Lemma~3.1]{IRW21} ---
and the $\TT{4}$-free tournaments on $4,5,6,7$ vertices number $3$, $3$,
$1$, $1$, the unique one on $7$ vertices being the Paley tournament
$QR_7$.
\end{proposition}

These inventories were computed by orbit-peeling of the labelled
enumerations under $S_k$, with a built-in closure assertion (every
relabelled image of every representative must reappear in the labelled
set), and re-derived by three independent enumerations in total; the
labelled orbit sizes cross-check exactly against the automorphism
groups (order $8$ for Bermond's graph, $21$ for $QR_7$).

Reversing all arcs preserves non-adjacency (hence $\I3$-freeness), maps
$\TT{4}$ to $\TT{4}$, and swaps $(p,q,s)\mapsto(q,p,s)$; both block
families of Proposition~\ref{prop:blocks} are closed under reversal. So
one may assume $p\le q$, and the possible types at $N=21$ are exactly
six. Assigning to each type every ordered choice of isomorphism classes
for the three blocks (both ordered pairs are kept when $p=q$) yields the
case list of Table~\ref{tab:types}: $346$ cases in all.

\begin{table}[ht]
\begin{center}
\begin{tabular}{cccc}
\toprule
type $(p,q,s)$ & block classes & cases \\
\midrule
$(5,8,7)$ & $25\times 1\times 1$ & $25$ \\
$(6,7,7)$ & $31\times 7\times 1$ & $217$ \\
$(6,8,6)$ & $31\times 1\times 1$ & $31$ \\
$(7,7,6)$ & $7\times 7\times 1$ & $49$ \\
$(7,8,5)$ & $7\times 1\times 3$ & $21$ \\
$(8,8,4)$ & $1\times 1\times 3$ & $3$ \\
\midrule
\multicolumn{2}{r}{total} & $346$ \\
\bottomrule
\end{tabular}
\end{center}
\caption{The six vertex types at $N=21$ (with $p\le q$) and the case
count.}
\label{tab:types}
\end{table}

\subsection{The encoding}\label{ssec:encoding}
The base formula $\mathrm{base}(N)$ has one Boolean variable per ordered
pair ($x_{u\to v}$: the arc $u\to v$ is present; $2\binom{N}{2}=420$
variables at $N=21$) and three clause families, transcribing the
definition and nothing else: for each pair, no $2$-cycle
($\lnot x_{u\to v}\lor\lnot x_{v\to u}$); for each triple, not
independent (the disjunction of its six arc variables); for each
$4$-set, for each of the $24$ transitive orientations, a clause
forbidding that orientation. At $N=21$ this is
$\binom{21}{2}+\binom{21}{3}+24\binom{21}{4}=145{,}180$ clauses. There
is no symmetry breaking of any kind in $\mathrm{base}(N)$: its models
correspond to free oriented graphs and to nothing else, and this
semantic claim was verified by brute force at small orders --- over all
$3^{\binom{N}{2}}$ orientations, the models of $\mathrm{base}(4)$ are
exactly the $612$ free graphs among $729$, and the models of
$\mathrm{base}(5)$ exactly the $35{,}488$ among $59{,}049$.

Each of the $346$ cases becomes a \emph{cube}: $\mathrm{base}(21)$
together with unit clauses fixing vertex $0$'s complete arc pattern
(arcs to $1,\dots,p$, arcs from $p+1,\dots,p+q$, non-adjacency to the
rest) and fixing each of the three blocks to its chosen class
representative --- $2\bigl(20+\binom{p}{2}+\binom{q}{2}+\binom{s}{2}\bigr)$
units in all, two per constrained pair. All arcs between different
blocks are left entirely free, constrained only by
$\mathrm{base}(21)$.

\subsection{Soundness of the case split}\label{ssec:canon}

\begin{proposition}\label{prop:canon}
Suppose some free oriented graph $D$ on $21$ vertices existed. Then, for
one of the $346$ cubes, some relabelling of $D$ (of $D$ with all arcs
reversed, if necessary) would satisfy every clause of that cube.
\end{proposition}

\begin{proof}
Fix any vertex $w$ of $D$ and let $(p,q,s)$ be its type; by
Lemma~\ref{lem:nbhd}, Lemma~\ref{lem:endpoints}, and $p+q+s=20$, the
type appears in Table~\ref{tab:types} after reversing all arcs if
$p>q$ (reversal preserves freeness and swaps the type; when $p=q$ both
ordered block pairs are among the cubes, so no choice is lost). By
Lemma~\ref{lem:nbhd} the three blocks induced by $\Np(w)$, $\Nm(w)$,
$\NI(w)$ belong to the enumerated families, and by the completeness of
the inventories (Proposition~\ref{prop:blocks}) each is isomorphic to
exactly one class representative. Relabel $w\mapsto 0$, $\Np(w)$ onto
$\{1,\dots,p\}$, $\Nm(w)$ onto $\{p+1,\dots,p+q\}$, and $\NI(w)$ onto
the rest, choosing within each block an isomorphism onto its
representative. The relabelled graph is free, so it satisfies
$\mathrm{base}(21)$ (whose models are exactly the free graphs), and by
construction it satisfies every unit clause of the cube indexed by its
type and block classes.
\end{proof}

Since all $346$ cubes are unsatisfiable --- \S\ref{ssec:certs} --- no
free graph on $21$ vertices exists, and with Theorem~\ref{thm:witness},
Theorem~\ref{thm:main} follows. The proof above is a short argument
whose computational inputs are Lemma~\ref{lem:endpoints}, the
completeness of Proposition~\ref{prop:blocks}, the semantics of
$\mathrm{base}(N)$, the identity of the $346$ cube files with
$\mathrm{base}(21)$-plus-units, and the $346$ refutations; every input
was machine-checked independently, most several times over
(\S\ref{sec:verif}).

\subsection{The certificates}\label{ssec:certs}
CaDiCaL~3.0.1 \cite{CaDiCaL} refuted each of the $346$ cubes,
emitting an LRAT proof \cite{LRAT} for each; every proof was verified by
an independently written checker that imports only the Python standard
library, accepts only on derivation of the empty clause, and rejects
deliberately corrupted controls. The verification ledger records, per
certificate, the verdict, the checked-step count, the SHA-256
digests of formula and of proof, and the proof's byte count. Across the
$N=21$ layer the certificates range from $236{,}431$ to $7{,}816{,}216$
checked steps, about $245$\,GB of uncompressed proof in total; the
complete corpus of the campaign (including the redundant layers of
\S\ref{ssec:redundant}) is $445$ certificates and just over $10^9$
checked steps. All $346$ solver logs report unsatisfiability; every
base has exactly one ledger line, verdict \textsc{verified}, formula
digest matching the audited file; no failures, no rejections.

\section{Verification and audit}\label{sec:verif}

\subsection{Completeness at content level}\label{ssec:complete}
The decomposition was audited by a referee-style completeness audit,
with independent code, and later re-derived once more from scratch: the
required cube inventory was recomputed from the type arithmetic and
fresh block enumerations, and each of the $346$ cube files was required
to equal $\mathrm{base}(21)$ plus its expected units \emph{as a clause
multiset} --- no missing clause, no duplicate, no stray unit, exact
DIMACS header --- with the match established one-to-one across the
whole layer, filename-blind. Negative controls (a dropped
clause, a duplicated clause, a hidden extra unit, a flipped literal)
were all rejected. The same audit passed for the $N=22$, $N=23$, and
$N=24$ layers.

As a positive control, the identical pipeline --- same generator, same
unit-fixing --- was run at $N=20$, where a free graph is known to exist:
the witness $W$, canonicalized per Proposition~\ref{prop:canon} onto its
class representatives, satisfies all $117{,}750$ clauses of the
generated $N=20$ cube; the solver reports satisfiability on both the
cube and the pure instance; and the extracted models were independently
re-verified free. A false-unsatisfiability pipeline would have failed
this control.

Two operational events are recorded for honesty. Mid-campaign the
checker's hint semantics were relaxed from strict format checking to
sound fixpoint semantics (skip hints unusable under the current
assignment, hard-fail on missing identifiers, accept only on the empty
clause) to accommodate a solver emission quirk; the change was ruled
sound, and all negative controls still fail. And three ledger lines were
poisoned by infrastructure failures and repaired: one when the checker
process was killed without output under memory and disk pressure, and
later two more (a second such kill, and a cloud-storage file eviction).
In every case only the verdict field changed --- never a formula and
never a proof. The two later originals survive in a ledger backup whose
formula and proof digests and byte counts are identical to the current
lines; the first survives verbatim in the failure log, the backup having
been taken afterwards. All three certificates verify with matching
digests, two of them under the referee's independent checker and one
regenerated byte-identically from its recorded solver command.

\subsection{The adversarial referee}\label{ssec:referee}
The confirmed result rests on a final adversarial referee pass, by a
separate agent instance under an explicit independence rule: every
load-bearing check used code written fresh in that session, sharing
nothing with the pipeline or the earlier audits beyond the DIMACS
variable-numbering convention. The referee re-derived the entire
skeleton from the definitions --- the neighbourhood lemma, both
endpoints of Lemma~\ref{lem:endpoints}, the inventories of
Proposition~\ref{prop:blocks}, the six types, the count $346$ --- and
re-established the content bijection between the $346$ required
formulas and the files on disk over regenerated content. It wrote and
control-validated its own LRAT checker, then attacked ten adversarially
chosen certificates: the largest proof in the corpus, all three entries
with a repair history (these four are the $2.0$--$2.3$\,GB proofs, four
of the five largest), the smallest, the largest of every remaining type,
and one further cube. All ten verified,
with step counts and digests equal to the ledger's in every field; in
all twelve cases where both checkers ran the same certificate, their
step counts agreed exactly. Eight of the ten proofs were regenerated
from nothing but the formula by re-running the recorded solver command;
all eight came out byte-identical to the recorded certificates. The
canonicalization of Proposition~\ref{prop:canon} was attacked with $57$
pivot canonicalizations across three graphs, including a family member
at $18$ vertices found by simulated annealing from a uniformly random
start, sharing no ancestry with the witness; no escape from the
inventory was found. Novelty was swept live the same day. On this
protocol the result was confirmed on August~11, 2026. The referee's
verdict record and fresh code are retained by the author and are not part
of the public deposit; what the deposit carries is the review log and the
separate verdict record for the $k(6,3)$ bound of \S\ref{ssec:k63}.

\section{Redundant layers, and $k(6,3)$}\label{sec:secondary}

\subsection{Certified exhaustions at $N=22,23,24$}\label{ssec:redundant}
Freeness is hereditary, so the $N=21$ exhaustion already implies that
no free graph exists on any larger order. The campaign nevertheless
produced certified exhaustions at
$N=23$ ($8$ cubes: types $(7,8,7)$ across the seven $7$-vertex block
classes, and $(8,8,6)$) and $N=22$ ($90$ cubes: types
$(6,8,7){\times}31$, $(7,7,7){\times}49$, $(7,8,6){\times}7$,
$(8,8,5){\times}3$), as well as a single structured instance at $N=24$;
all $99$ certificates are verified and ledgered, and the multiset
completeness audit of \S\ref{ssec:complete} passed for each layer.
These layers are logically redundant given Theorem~\ref{thm:main}; they
stand as independent confirmations along the descent
$25\to 24\to 23\to 22\to 21$, and their certificates are part of the
corpus.

\subsection{The bracket for $k(6,3)$}\label{ssec:k63}

\begin{proof}[Proof of Proposition~\ref{prop:k63}]
Upper bound: $k(m,3)\le m^2-m+3$ is \cite[Proposition~3.4]{IRW21};
at $m=6$ this is $33$. Lower bound: let $S=\{3,8,10,12,17\}\subset
\mathbb{Z}_{28}$ and let $D=\mathrm{Cay}(\mathbb{Z}_{28},S)$, with arcs
$x\to x+s$ for $s\in S$. Since $-S=\{11,16,18,20,25\}$ is disjoint from
$S$, $D$ is an oriented graph; it is $5$-regular in and out, with $140$
arcs. A transitive triangle $x\to y\to z$, $x\to z$ in
$\mathrm{Cay}(G,S)$ is exactly a pair $a=y-x$, $b=z-y$ in $S$ with
$a+b=z-x$ also in $S$; here
$S+S=\{1,6,11,13,15,16,18,20,22,24,25,27\}$ is
disjoint from $S$, so $D$ has no $\TT{3}$. The underlying graph of $D$
is $\mathrm{Cay}(\mathbb{Z}_{28},S\cup-S)$, whose independence number is
$5$; so $D$ has no $\I6$. Hence $k(6,3)\ge 29$.
\end{proof}

Only the independence number is a machine check, over the
$\binom{28}{6}=376{,}740$ six-subsets. It was performed by the
independently written witness verifiers of \S\ref{sec:witness}, and
again by exact computation of the independence number by branch and
bound. Two further $28$-vertex
witnesses, one of them a Cayley digraph over
$\mathbb{Z}_{14}\times\mathbb{Z}_2$, verify equally and are archived.

Disjoint unions cannot reach $28$: splitting the independence budget
$5$ across components with the extremal component orders available
($3$ at budget $1$, since a $\TT3$-free tournament has at most $3$
vertices; $k(a{+}1,3)-1$ in general) gives $3+22=25$ as the best
partition, so every witness on $26$ or more vertices --- including
$D$, whose connection set contains a generator --- is connected in the
underlying adjacency. No lower bound beyond $24$ ($W_{22}$, the
$\{\I5,\TT{3}\}$-free circulant of \cite[Observation~4.2]{IRW21}, plus
an isolated vertex) appears in the literature.

\section{Certification}\label{sec:certification}

The permanent record of the computation is deliberately small: the
verification ledger (one line per certificate: verdict, checked-step
count, SHA-256 of formula and of proof, and the proof's byte count), the
queue files recording the exact solver command line for every cube of
the $N=21$, $22$ and $23$ layers, the generator and checker scripts, the
block-inventory representatives the generator consumes, the witness
files, and the audit and referee
records. That record is split between a public deposit and a private
regeneration kit.

Deposited with this note in the program's public certificate
repository \emph{Certify} (\url{https://github.com/05oz/certify};
concept DOI \href{https://doi.org/10.5281/zenodo.21799111}%
{10.5281/zenodo.21799111}), version DOI
\href{https://doi.org/10.5281/zenodo.21890619}{10.5281/zenodo.21890619},
are the ledger, the regeneration instructions, the two witness files,
six scripts --- the encoder \texttt{gen\_cnf.py}, the structured-cube
generator \texttt{make\_structured.py}, the proof checker
\texttt{lrat\_check.py}, the witness checker
\texttt{verify\_witness.py}, and the audits \texttt{audit\_cnf.py} and
\texttt{audit\_multiset.py} --- and the review log together with the
referee verdict for $k(6,3)$. Retained by the author and not
redistributed are the queue files, the $70$ block-class representatives
the generators consume, the rest of the checking and audit pipeline, and
the audit and $k(3,4)$ referee records. So the lower-bound half replays
from the deposit outright, as does any certificate whose formula and
proof the reader has; regenerating a case formula to its recorded
\texttt{cnf\_sha256} needs the representatives, and
\texttt{audit\_multiset.py} imports a module of the kit and does not run
from the deposit alone.

The multi-gigabyte proofs themselves are a cache, not the record: each
regenerates from its recorded command --- deterministically, in our
environment: the referee's eight regenerations were all byte-identical
to the originals, matching the ledger's digests --- and the deposit
documents the regeneration procedure and what it requires. As to
tooling, the checkers are standard-library only, so replaying a
certificate or a witness requires only CPython; regenerating proofs
requires CaDiCaL~3.0.1.

\section{The trusted base}\label{sec:tcb}

What a skeptic must believe, separated from what they can replay.

\emph{Machine-checked, replayable with standard-library Python:} the
witness $W$ and the $28$-vertex witness of
Proposition~\ref{prop:k63} (milliseconds); the labelled enumerations
behind Lemma~\ref{lem:endpoints} and Proposition~\ref{prop:blocks},
performed three independent times, with closure assertions (minutes);
the brute-force semantics of $\mathrm{base}(4)$ and $\mathrm{base}(5)$
over all orientations; the clause-multiset identity of all $445$ cube
files with their re-derived specifications (this audit runs in the
regeneration kit, not from the deposit alone); and the LRAT replay of all
$445$ certificates by an independently written checker, with twelve of
them re-run by a second, unrelated checker agreeing on the exact step
count in every case. Both checkers reject all negative controls.

\emph{Trusted:}
\begin{enumerate}
\item \emph{The checkers' proof semantics.} Both checkers implement
RUP-with-hints with sound fixpoint treatment of hints and accept only on
derivation of the empty clause. That this discipline is sound is a short
argument on paper, not a machine-checked one; the mitigation is two
independent implementations in agreement, plus negative controls
(corrupted proof, mutated formula, dropped hint) rejected by both.
\item \emph{The assembly of \S\ref{ssec:canon}.} The proof of
Proposition~\ref{prop:canon} is half a page of mathematics to be read by
a human. Every computational input to it is machine-checked; the gluing
is not, and is short by design.
\item \emph{The generalization of the base semantics from $N=4,5$ to
$N=21$.} The brute-force model check is exhaustive only at $N=4,5$; at
$N=21$ the claim rests on the clause families being uniform in $N$
(with closed-form clause count
$\binom{N}{2}+\binom{N}{3}+24\binom{N}{4}$ matching at every audited
order) and on the independent re-derivation of the exact clause
multiset from the definitions.
\item \emph{SHA-256 collision resistance}, for the bookkeeping that ties
ledger lines to files. A skeptic can replay every shipped LRAT certificate
without trusting any hash; regenerating a case formula to its recorded
\texttt{cnf\_sha256}, however, needs the block representatives, which are
part of the private regeneration kit and not of this deposit.
\item \emph{CPython, the operating system, and the hardware.}
\end{enumerate}

\noindent CaDiCaL, the cube generators, and the solving pipeline are
\emph{not} in the trusted base: nothing depends on the solver being
correct, only on its proofs checking, and every proof was checked by
code unrelated to the solver --- a sample of them by two such checkers
independently.

\section{Questions}\label{sec:questions}

What is this computation \emph{for}? Four directions it opens.

\subsection{The extremal structure at $(3,4)$}\label{ssec:qstruct}
The extremal graph for $k(3,3)$ is unique up to isomorphism
\cite[Lemma~3.1]{IRW21} --- a circulant. The witness $W$ of
\S\ref{sec:witness} shows no
such symmetry: it has four vertex types, so it is not vertex-transitive,
although the Paley tournament $QR_7$ appears as the non-neighbourhood of
eleven of its twenty vertices. This mirrors the situation at $k(4,3)$,
where \cite{IRW21} remark that no Cayley witness on $14$ vertices
exists, while the $k(5,3)$ witness is again circulant.

\begin{question}\label{q:unique}
Is $W$ the unique $\{\I3,\TT{4}\}$-free oriented graph on $20$ vertices
up to isomorphism? If not, what is the number of extremal graphs, and
does any admit an algebraic description?
\end{question}

A related quantitative puzzle: the counting bound of
\S\ref{ssec:decomp} caps free graphs at $24$ vertices, yet the largest
free graph has $20$ --- at $N=21,22,23$ every one of the $346+90+8$
cases dies, with no tightness anywhere. A structural explanation of
this slack of $4$ (at $k(3,3)$ the analogous cap $1+2+2+3=8$ is
attained, by Bermond's graph) is exactly what a human proof would have
to supply.

\subsection{The gap at $k(6,3)$}\label{ssec:qk63}
Proposition~\ref{prop:k63} leaves $k(6,3)\in[29,33]$, now the smallest
open case of the $k(m,3)$ column. Every witness on $26$ or more vertices
is connected in the underlying adjacency, which points toward
vertex-transitive candidates: for Cayley digraphs, $\TT{3}$-freeness of
$\mathrm{Cay}(G,S)$ is exactly sum-freeness of $S$, which makes the
search cheap. This is how the $28$-vertex witness was found. An
exhaustive sweep of connection sets over the cyclic groups of orders
$29$, $30$ and $31$ --- the only abelian groups of those orders --- and
over three of the seven abelian groups of order $32$ produced nothing,
so any witness beating $28$ lies outside that range of constructions. The
upper bound looks harder: an exhaustion in this note's style
would need complete inventories of $\{\I5,\TT{3}\}$-free graphs on up
to $22$ vertices, far beyond present enumeration, so progress below
$33$ seems likelier to come from sharpening the $m^2-m+3$ argument of
\cite{IRW21}.

\subsection{$k(4,4)$ and beyond}\label{ssec:qk44}
The new value propagates through the standard recursion. By
Lemma~\ref{lem:nbhd} in its general form \cite[Lemma~2.1]{IRW21},
every $\{\I4,\TT{4}\}$-free graph has $d^{\pm}(v)\le k(4,3)-1=14$ and
$|\NI(v)|\le k(3,4)-1=20$, so such a graph has at most
$1+14+14+20=49$ vertices, and likewise every $\{\I3,\TT{5}\}$-free
graph has at most $1+20+20+13=54$; hence
\[
k(4,4)\;\le\;50,
\qquad\qquad
k(3,5)\;\le\;55 ,
\]
which is the recursion of \cite[Lemma~2.3]{IRW21} evaluated at the new
value. Meanwhile $k(4,4)\ge k(3,4)=21$, since a graph with no $\I3$ has
no $\I4$,
so $W$ is also $\{\I4,\TT{4}\}$-free. An honest feasibility assessment:
determining $k(4,4)$ or $k(3,5)$ by the present method is out of reach.
The decomposition would require complete isomorph inventories of
$\{\I4,\TT{3}\}$-free graphs up to $14$ vertices, or of
$\{\I3,\TT{4}\}$-free graphs up to $20$ vertices --- and the labelled
counts of the latter family are already $35{,}488$ at $5$ vertices and
$4{,}490{,}572$ at $6$ --- and the resulting instances at $N$ in the
forties would dwarf the $10^9$ checked steps spent here at $N=21$. What
does look feasible: raising lower bounds by witness search (any
verified witness is a one-line check), and isomorph-free generation of
the free families at small orders to map the terrain.

\subsection{Toward a human-readable proof}\label{ssec:qhuman}
The upper-bound half of Theorem~\ref{thm:main} is $346$ machine proofs
averaging about two and a half million checked steps. Two concrete paths would shrink
it. First, structure theory: the proof
of Proposition~\ref{prop:canon} shows that any free graph on $21$
vertices is assembled from three blocks drawn from very short lists
around every vertex simultaneously; a human argument that such an
assembly forces an $\I3$ or a $\TT{4}$ --- perhaps via the $QR_7$ and
Bermond blocks, which dominate the type table --- would replace the
certificates outright, and would likely explain the slack of
\S\ref{ssec:qstruct} as a by-product. Second, proof compression: the
LRAT corpus is an upper bound on the difficulty of the theorem, not a
measure of it, and nothing is known about the minimum certificate size;
a decomposition chosen for proof economy rather than for engineering
convenience might cut the corpus by orders of magnitude and expose
which cases are genuinely hard. Either advance would turn a determined
value into an understood one.

\section*{Acknowledgments}

The problem is Erd\H{o}s and Rado's; its continued visibility owes much
to Thomas Bloom's Erd\H{o}s problems collection \cite{Blo}. The line of
exact values this note extends was opened by Jean-Claude Bermond and
carried to $k(4,3)$ and $k(5,3)$ by Ferdinand Ihringer, Deepak
Rajendraprasad, and Thilo Weinert, whose paper also named this note's
target. The certificate infrastructure rests on CaDiCaL by Armin Biere
and coauthors, on the LRAT format of Cruz-Filipe, Heule, Hunt, Kaufmann,
and Schneider-Kamp, and on the cube-and-conquer paradigm of Heule,
Kullmann, Wieringa, and Biere. The computation and drafting were
AI-assisted as stated in the first footnote; the adversarial referee
protocol of \S\ref{ssec:referee} and the referee's rule that a
computation-only paper close with the questions it opens shaped this
note's final form.

\begin{thebibliography}{99}
\footnotesize

\bibitem[Ber74]{Ber74}
J.-C. Bermond,
\emph{Some Ramsey numbers for directed graphs},
Discrete Math.\ \textbf{9} (1974), 313--321.

\bibitem[BFF{+}24]{CaDiCaL}
A.~Biere, T.~Faller, K.~Fazekas, M.~Fleury, N.~Froleyks and F.~Pollitt,
\emph{CaDiCaL 2.0},
in: Computer Aided Verification --- CAV 2024, Lecture Notes in Comput.\
Sci.\ \textbf{14681}, Springer, 2024, 133--152. Version used to produce
the proofs: CaDiCaL 3.0.1 (\texttt{--lrat}).

\bibitem[Blo]{Blo}
T.~F. Bloom, \emph{Erd\H{o}s Problem \#112},
\url{https://www.erdosproblems.com/112}, accessed 2026-08-11.

\bibitem[CFHKS17]{LRAT}
L.~Cruz-Filipe, M.~J.~H. Heule, W.~A. Hunt~Jr., M.~Kaufmann and
P.~Schneider-Kamp,
\emph{Efficient certified RAT verification},
in: Automated Deduction --- CADE-26, Lecture Notes in Comput.\ Sci.\
\textbf{10395}, Springer, 2017, 220--236.

\bibitem[ErRa67]{ErRa67}
P.~Erd\H{o}s and R.~Rado,
\emph{Partition relations and transitivity domains of binary relations},
J.\ London Math.\ Soc.\ \textbf{42} (1967), 624--633.

\bibitem[HKM16]{HKM16}
M.~J.~H. Heule, O.~Kullmann and V.~W. Marek,
\emph{Solving and verifying the Boolean Pythagorean triples problem via
cube-and-conquer},
in: Theory and Applications of Satisfiability Testing --- SAT 2016,
Lecture Notes in Comput.\ Sci.\ \textbf{9710}, Springer, 2016, 228--245.

\bibitem[HKWB11]{HKWB11}
M.~J.~H. Heule, O.~Kullmann, S.~Wieringa and A.~Biere,
\emph{Cube and conquer: guiding CDCL SAT solvers by lookaheads},
in: Hardware and Software: Verification and Testing --- HVC 2011,
Lecture Notes in Comput.\ Sci.\ \textbf{7261}, Springer, 2012, 50--65.

\bibitem[IRW21]{IRW21}
F.~Ihringer, D.~Rajendraprasad and T.~Weinert,
\emph{New bounds on the Ramsey number $r(I_m,L_n)$},
Discrete Math.\ \textbf{344} (2021), no.~3, 112268. Also arXiv:1707.09556 (v3,
April 2020), the text read here (see the second footnote).

\bibitem[LaMi97]{LaMi97}
J.~A. Larson and W.~J. Mitchell,
\emph{On a problem of Erd\H{o}s and Rado},
Ann.\ Comb.\ \textbf{1} (1997), 245--252.

\end{thebibliography}

\end{document}
